\documentclass[10pt]{article}
\usepackage[T1]{fontenc}
\usepackage[utf8]{inputenc}
\usepackage[french]{babel}
\usepackage{amsmath,amssymb,mathrsfs,mathpazo}
\usepackage{geometry}
\geometry{a4paper,margin=2cm}
\usepackage{graphicx,array,tabularx,tabularray,colortbl,multirow,multicol,booktabs,enumitem,pifont,ulem,fancybox,diagbox,cancel}
\usepackage{pstricks,pstricks-add,pst-plot,pst-tree,pst-node,pst-text,pst-func,pst-math,pst-bspline,pst-3dplot,pst-eucl,pst-solides3d}
\let\pscalculate\undefined
\usepackage{pst-bezier}
\usepackage{tikz,pgfplots,tkz-tab}
\usetikzlibrary{arrows,arrows.meta,patterns,calc,positioning,trees,decorations.markings,shapes.geometric,shapes.misc,shapes.arrows,intersections,decorations.pathmorphing,decorations.pathreplacing,backgrounds,shadows,babel}
\pgfplotsset{compat=1.18}
\usepgfplotslibrary{fillbetween}
\usepackage[np]{numprint}
\usepackage{listings,makecell,pas-tableur,twoopt}
\let\barre\undefined
\usepackage[pstricks]{tablvar}
\definecolor{darkgreen}{RGB}{0,100,0}
\providecommand{\red}{\color{red}}
\providecommand{\blue}{\color{blue}}
\providecommand{\green}{\color{green}}
\providecommand{\black}{\color{black}}
\providecommand{\white}{\color{white}}
\providecommand{\gray}{\color{gray}}
\providecommand{\cyan}{\color{cyan}}
\providecommand{\magenta}{\color{magenta}}
\providecommand{\yellow}{\color{yellow}}
\providecommand{\violet}{\color{violet}}
\providecommand{\orange}{\color{orange}}
\providecommand{\euro}{\texteuro}
\providecommand{\EUR}{\texteuro}
\providecommand{\ieme}{\textsuperscript{e}}
\providecommand{\degres}{\ensuremath{^\circ}}
\providecommand{\degree}{\ensuremath{^\circ}}
\providecommand{\vv}[1]{\overrightarrow{#1}}
\providecommand{\Coord}[1]{#1}
\providecommand{\decosix}{}
\providecommand{\starredbullet}{\ensuremath{\bullet}}
\pagestyle{empty}
\setlength{\parindent}{0pt}
\usepackage{adjustbox}\DeclareUnicodeCharacter{2236}{:}\DeclareUnicodeCharacter{2190}{\ensuremath{\leftarrow}}
\providecommand{\R}{\mathbb{R}}
\providecommand{\N}{\mathbb{N}}
\providecommand{\Z}{\mathbb{Z}}
\providecommand{\Q}{\mathbb{Q}}
\providecommand{\D}{\mathbb{D}}
\providecommand{\C}{\mathbb{C}}
\providecommand{\e}{{\mathrm{e}}}
\providecommand{\vect}[1]{\overrightarrow{#1}}
\providecommand{\vectt}[1]{\overrightarrow{\mathrm{#1}}}
\providecommand{\vvt}[1]{\overrightarrow{\mathrm{#1}}}
\providecommand{\np}[2][]{{#2~#1}}
\providecommand{\vv}[1]{{\overrightarrow{#1}}}
\renewcommand{\d}{{\mathrm{d}}}
\providecommand{\degres}{{^\circ}}
\providecommand{\degre}{{^\circ}}
\providecommand{\blue}{\color{blue}}
\providecommand{\red}{\color{red}}
\providecommand{\enskip}{\quad}
\providecommand{\llbracket}{\lbrack\!\lbrack}
\providecommand{\rrbracket}{\rbrack\!\rbrack}
\providecommand{\euro}{{\text{€}}}
\providecommand{\no}{{\text{n°}}}
\providecommand{\interval}[2]{{[#1;#2]}}
\providecommand{\intervFO}[2]{{[#1;#2[}}
\providecommand{\dotfill}{\ldots}

\providecommand{\ds}{\displaystyle}
\providecommand{\textding}[1]{#1}
\providecommand{\Oij}{\(\mathrm{O};\vec{\imath},\vec{\jmath}\)}
\providecommand{\Oijk}{\(\mathrm{O};\vec{\imath},\vec{\jmath},\vec{k}\)}
\providecommand{\Ouv}{\(\mathrm{O};\vec{u},\vec{v}\)}
\providecommand{\intervFF}[2]{\ensuremath{[#1;#2]}}
\renewcommand{\euro}{\ifmmode\text{\texteuro}\else\texteuro\fi}
\DeclareUnicodeCharacter{00B0}{\ifmmode^{\circ}\else\textdegree\fi}
\DeclareUnicodeCharacter{20AC}{\ifmmode\text{\texteuro}\else\texteuro\fi}
% Énoncés : Éducation nationale. Transcription et corrigés : APMEP, Denis Vergès et contributeurs. Découpage et classement : CoopMaths, Sébastien Lozano, Éric Elter et contributeurs. Adaptation de la présentation : Maths Explorer.
% CC BY-SA : https://coopmaths.fr/www/about/
% https://coopmaths.fr/alea/static/bac/2023/tex/bac_2023_08_sujet1_caledonie_1.tex
\begin{document}
\noindent\textit{APMEP / CoopMaths. CC BY-SA. Présentation adaptée par Maths Explorer.}\par\medskip
\medskip

Une entreprise de location de bateaux de tourisme propose à ses clients deux types de bateaux : bateau à voile et bateau à moteur.

Par ailleurs, un client peut prendre l'option PILOTE. Dans ce cas, le bateau, qu'il soit à voile ou à moteur, est loué avec un pilote.

On sait que:
\begin{itemize}
\item[$\bullet~$] 60\,\% des clients choisissent un bateau à voile;
parmi eux, 20\,\% prennent l'option PILOTE.
\item[$\bullet~$] 42\,\% des clients prennent l'option PILOTE.
\end{itemize}

On choisit au hasard un client et on considère les évènements :

\begin{itemize}
\item[$\bullet~$]$V$: \og le client choisit un bateau à voile\fg{} ;
\item[$\bullet~$]$L$: \og le client prend l'option PILOTE ».
\end{itemize}

\begin{center}\emph{Les trois parties peuvent être traitées de manière indépendante}\end{center}

\medskip

\textbf{Partie A}

\medskip

\medskip

\begin{enumerate}
\item Traduire la situation par un arbre pondéré que l'on complètera au fur et à mesure.
\item Calculer la probabilité que le client choisisse un bateau à voile et qu'il ne prenne pas l'option PILOTE.
\item Démontrer que la probabilité que le client choisisse un bateau à moteur et qu'il prenne l'option PILOTE est égale à $0,30$.
\item En déduire $P_{\overline{V}}(L)$, probabilité de $L$ sachant que $V$ n'est pas réalisé.
\item Un client a pris l'option PILOTE.

Quelle est la probabilité qu'il ait choisi un bateau à voile ? Arrondir à $0,01$ près.
\end{enumerate}

\medskip

\textbf{Partie B}

\medskip
Lorsqu'un client ne prend pas l'option PILOTE, la probabilité que son bateau subisse une avarie est égale à $0,12$. Cette probabilité n'est que de $0,005$ si le client prend l'option PILOTE.

On considère un client. On note $A$ l'évènement: \og son bateau subit une avarie \fg.

\medskip

\begin{enumerate}
\item Déterminer $P(L \cap A)$ et $P\left(\overline{L} \cap A\right)$.
\item L'entreprise loue \np{1000} bateaux.

À combien d'avaries peut-elle s'attendre ?
\end{enumerate}

\medskip

\textbf{Partie C}

\medskip
On rappelle que la probabilité qu'un client donné prenne l'option PILOTE est égale à 0,42.

On considère un échantillon aléatoire de $40$ clients. On note $X$ la variable aléatoire comptant le nombre de clients de l'échantillon prenant l'option PILOTE.

\medskip

\begin{enumerate}
\item On admet que la variable aléatoire $X$ suit une loi binomiale. Donner sans justification ses paramètres.
\item Calculer la probabilité, arrondie à $10^{-3}$, qu'au moins $1$5~clients prennent l'option PILOTE.
\end{enumerate}

\bigskip

\end{document}
